% Part: incompleteness
% Chapter: theories-computability
% Section: oconsis-ext-of-q-undec

\documentclass[../../../include/open-logic-section]{subfiles}

\begin{document}

\olfileid[fr]{inc}{tcp}{oqn}

\olsection{Les extensions \texorpdfstring{$\omega$}{ω}-cohérentes de $\Th{Q}$ sont indécidables}

\begin{explain}
La preuve que l'ensemble des théorèmes de~$\Th{Q}$ est
complet parmi les ensembles effectivement énumérables
utilisait le fait que tout !!{sentence} démontrable
dans~$\Th{Q}$ est « vrai » dans les entiers naturels.
La définition et le théorème suivants renforcent ce
résultat en isolant la seule propriété de cette vérité
dont la preuve avait besoin. Il n'est pourtant pas
nécessaire de s'y attarder : nous obtiendrons bientôt
un résultat plus fort. Nous présentons surtout celui-ci
pour des raisons historiques. L'article original de
Gödel employait la $\omega$-cohérence, mais son résultat
a rapidement été renforcé en ne supposant que la
cohérence ordinaire.
\end{explain}

\begin{defn}

Une théorie~$\Th{T}$ est \emph{$\omega$-cohérente} si,
pour tout !!{sentence} de la forme
$\lexists[x][!A(x)]$, lorsque $\Th{T}$ démontre
$\lnot !A(\num 0)$, $\lnot !A(\num 1)$,
$\lnot !A(\num 2)$, \dots, elle ne démontre pas
$\lexists[x][!A(x)]$.
\end{defn}

\begin{thm}\ollabel{thm:oconsis-q}
Toute théorie~$\Th{T}$ qui contient~$\Th{Q}$ et qui
est $\omega$-cohérente est indécidable.\footnote{Dans la source, le renvoi à ce théorème est attaché à la définition précédente; sa destination a été corrigée.}
\end{thm}

\begin{proof}
Puisque $\Th{T}$ contient~$\Th{Q}$, elle représente
les fonctions et les relations calculables.
Reprenons la réduction précédente. Si $x\in K$,
$\Th{T}$ démontre
$\lexists[s][!A_T(\num x,\num x,s)]$.
Réciproquement, supposons qu'elle démontre cet
!!{sentence}. Si $x\notin K$, la machine de code~$x$
ne s'arrête pas sur l'entrée~$x$. Comme $!A_T$
représente la relation de Kleene~$T$,
$\Th{Q}$, et donc $\Th{T}$, démontre chacune
des !!{sentence}s
$\lnot !A_T(\num x,\num x,\num 0)$,
$\lnot !A_T(\num x,\num x,\num 1)$, \dots.
En démontrant aussi l'existentielle, $\Th{T}$
serait $\omega$-incohérente. Donc $x\in K$
si et seulement si $\Th{T}$ démontre
$\lexists[s][!A_T(\num x,\num x,s)]$.
Le code de cet !!{sentence} se calcule à partir
de~$x$; si $\Th{T}$ était décidable, $K$ le
serait aussi, ce qui est impossible.
\end{proof}

\end{document}
